\chapter{Curvas Simples. Clasificación. Teorema de la Curva de Jordan. Teorema de la Divergencia}

\begin{definicion}
    Una \textbf{curva simple} es un subconjunto $C\subset \bb{R}^2$ no vacío tal que $\forall p\in C$ se tiene que 
    \begin{itemize}
        \item $\exists\ V$ entorno abierto de $p$ en $C$.
        \item $\exists\ I\subset \bb{R}$ intervalo abierto.
        \item $\exists\ \alpha:I \to \bb{R}^2$ curva parametrizada regular
    \end{itemize} 
    de forma que $\alpha(I)=V$, y la aplicación $\alpha:I\to V$ es un homeomorfismo.
\end{definicion}

% dibujo

% En cierta forma estamos copiando la definición de superficies pero para R^2.

\begin{observacion}\
    \begin{enumerate}
        \item La aplicación $\alpha$ la llamaremos \textbf{parametrización local} de $C$.
        \item $C$ no puede tener ``esquinas'' (o ``picos'') ya que $\alpha$ es regular.
        \item $C$ no puede tener autointersecciones. Ni siquiera el Folium de Descartes puede ser una curva simple, a pesar de que en este caso $\alpha$ es inyectiva.
        \item Los entornos $V$ son todos de la forma
        \begin{gather*}
            V = U\cap C
        \end{gather*}
        donde $U$ es un entorno abierto de $p$ en $\bb{R}^2$.
        \item Los entornos $V$ son todos conexos.
    \end{enumerate}
\end{observacion}

\begin{ejemplo}
    Supongamos que tenemos una curva simple $C\subset \bb{R}^2$ y consideramos las parametrizaciones $\alpha: I \to C$ y $\beta:J \to C$ de forma que $W = \alpha(I)\cap \beta(J)\neq \emptyset$.\\

    Tenemos que $\alpha^{-1}(W)\subset I$ y que $\beta^{-1}(W)=J$. Además, como $W$ es abiertos, tendremos que $\alpha^{-1}(W)$ es abierto en $I$ y $\beta^{-1}(W)$ es abierto en $J$ y por ser abiertos de un abierto, ambos serán abiertos en $\bb{R}$.\\

    Además, podemos considerar los \textbf{cambios de parametrización}
    \begin{align*}
        \phi &= \beta^{-1}\circ \alpha : \alpha^{-1}(W) \to \beta^{-1}(W)\\
        \phi^{-1} = \psi &= \alpha^{-1}\circ \beta : \beta^{-1}(W) \to \circ^{-1}(W)
    \end{align*}
    Veamos que $\phi$ y $\psi$ son \textbf{difeomorfismos}. Para ello, tan solo tendremos que verlo para $\phi$. Además, bastará con ver que $\phi$ es derivable (diferenciable), ya que tenemos que es una biyección y su inversa será $\psi$, donde $\alpha$ y $\beta$ juegan papeles inversos por lo que será equivalente a ver que $\phi$ lo sea.\\

    Recordamos la situación en la que estamos
    \begin{gather*}
        \xymatrix@R=3mm{
            \alpha^{-1}(W) \ar[r]^{\phi}  & \beta^{-1}(W) \ar[r]^{\beta_1} & \bb{R} \\
            t_0 \ar@{|->}[r] & \phi(t_0) \ar@{|->}[r] & \beta(\phi(t_0))
        }
    \end{gather*}
    Además, como $\beta$ es regular tenemos que $\beta'(\phi(t_0)) \neq (0,0)$. Supongamos sin pérdida de generalidad que $\beta_1(\phi(t_0)) \neq 0$ y estamos en condiciones de aplicarle el Teorema de la Función Inversa. Este nos dice que existe un $W_2$ entorno abierto de $\phi(t_0)$ en $\beta^{-1}(W)$ (que también será abierto en $\bb{R}$ por ser abierto de abierto) y existe también un $W_1$ entorno abierto de $\beta_1(\phi(t_0))$ en $\bb{R}$ de forma que $\beta_1(W_2) = W_1$ y tal que $\beta_{1|_{W_2}} : W_2 \to W_1$ es un difeomorfismo. Además,
    \begin{gather*}
        \phi = \beta^{-1}\circ \alpha \quad \Rightarrow \quad \alpha = \beta \circ \phi \quad \Rightarrow \quad \alpha_1 = \beta_1 \circ \phi
    \end{gather*}
    Como $\phi$ es continua, en particular será continua en $t_0$ por lo que tenemos un entorno abierto de $\phi(t_0)$ que es $W_2$. Por la definición topológica de continuidad tendremos que existe un $W_3$ entorno abierto de $t_0$ en $\alpha^{-1}(W)$ y por tanto en $\bb{R}$ de forma que $\phi(W_3) \subset W_2$.\\

    Vamos ahora a considerar la restricción
    \begin{gather*}
        \alpha_{1|_{W_3}} = (\beta_1\circ \phi)_{|_{W_3}} = \beta_{1|_{W_2}} \circ \phi_{|_{W_3}}
    \end{gather*}
    Además, por ser una parametrización, $\alpha$ es derivable y por tanto su primera componente $\alpha_1$ también lo será (y en particular su restricción también). Del Teorema de la función inversa anterior sabemos que $\beta_{1|_{W_2}}$ es un difeomorfismo. Por tanto podemos pasar 
    \begin{gather*}
        \phi_{|_{W_3}} = \beta_{1|_{W_2}}^{-1} \circ \alpha_{1|_{W_3}}
    \end{gather*}
    y llegamos a la conclusión de que $\phi_{|_{W_3}}$ es derivable, pero $W_3$ es un entorno abierto de $t_0$. Hemos llegado a que para un $t_0$ cualquiera del dominio del cambio de parámetros existe un entorno suyo en el que $\phi$ es derivable. En particular lo será en $t_0$.
\end{ejemplo}

% No hay clase el 17 de septiembre ni del 28-1.
% Se recuperará:
%   - día 25 de septiembre de 12:00-14:00.
%   - día 9 de octubre de 12:00-14:00.
%   - día 16 de octubre de 13:00-14:00.


Del ejemplo anterior podemos concluir la siguiente proposición:\\
\begin{prop}
    Sean $C\subset \bb{R}^2$ una curva simple y $\alpha:I \to C$, $\beta: J \to C$ dos parametrizaciones suyas cuyas imágenes se solapan, es decir, tal que 
    \begin{gather*}
        W = \alpha(I)\cap \beta(J) \neq \emptyset
    \end{gather*}
    Entonces las aplicaciones \textbf{cambio de parámetros} siguientes
    \begin{align*}
        \phi &= \beta^{-1}\circ \alpha : \alpha^{-1}(W) \to \beta^{-1}(W)\\
        \psi &= \alpha^{-1}\circ \beta : \beta^{-1}(W) \to \circ^{-1}(W)
    \end{align*}
    son \underline{difeomorfismos} entre abiertos de $\bb{R}$ (de $I$ y $J$ respectivamente).
\end{prop}

\begin{definicion}
    Sea $C\subset \bb{R}^2$ una curva simple y un punto $p\in C$. Definimos la \textbf{recta tangente} a la curva $C$ en el punto $p$ como la recta afín $T_pC$ que pasa por $p$ y tiene dirección $\alpha'(t_0)$, es decir,
    \begin{gather*}
        T_pC = p + \langle \alpha'(t_0) \rangle
    \end{gather*}
    donde $\alpha: I \to C$ es una parametrización con $\alpha(t_0)=p$.
\end{definicion}

\begin{observacion}
    En primer lugar es fácil ver que $T_pC$ es una recta, ya que $\alpha'(t_0)\neq 0$ por ser una curva regular. \\
    
    Nos planteamos además, qué ocurre si tomamos una parametrización distinta. Para ello consideramos $\beta: J \to C$ otra parametrización con $\beta(u_0)=p$, donde $u_0\in J$ y queremos ver la relación entre $\alpha'(t_0)$ y $\beta(u_0)$.\\

    Sabemos que $\phi = \beta^{-1} \circ \alpha$ es un difeomorfismo. Además, esto podemos aplicarlo porque sabemos que se solapan, ya que 
    \begin{gather*}
        \left.
        \begin{array}{l}
            p = \alpha(t_0) \in \alpha (I)\\
            p = \beta(u_0) \in \beta(J)
        \end{array}
        \right\} \Rightarrow \alpha(I) \cap \beta(J) \neq \emptyset,\quad \phi(t_0) = u_0
    \end{gather*} 
    Además $\alpha = \beta \circ \phi$ y estamos en condiciones de derivar en $t_0$ aplicando la regla de la cadena.
    \begin{gather*}
        \alpha'(t_0) = (\beta \circ \phi)'(t_0) = \phi'(t_0)\beta'(\phi(t_0)) = \phi'(t_0)\beta'(u_0) \neq 0
    \end{gather*}
    Por colinealidad tenemos que 
    \begin{gather*}
        \langle \alpha'(t_0) \rangle = \langle \beta'(t_0) \rangle
    \end{gather*}
    lo que nos confirma que la definición anterior está bien hecha y no depende de la parametrización escogida.
\end{observacion}

\begin{ejemplo}\
    \begin{enumerate}
        \item Sea $C\subset \bb{R}^2$ una curva simple y consideramos $C'\subset C$ n abierto suyo. Entonces $C'$ es también una curva simple.
        \begin{proof}
            Para probarlo tenemos que ver que para todo punto existe una parametrización regular. Así consideramos un $p'\in C'$ y en particular $p'\in C$ por lo que existirá un $V\subset C$ entorno de $p'$, un intervalo abierto $I\subset \bb{R}$ y una curva regular $\alpha : I \to \bb{R}^2$ tal que $\alpha(I)=V$ y tal que $\alpha : I \to V$ es un homeomorfismo.\\

            Podríamos ahora pensar en elegir un $V'=V\cap C'$ pero a poco que lo pensemos no vale porque $V'$ puede ser no conexo. Por ello consideramos $V'$ como la componente conexa de $V\cap C'$ que contiene a $p'$, que será un entorno abierto y conexo de $p'$ en $C'$.\\

            Con esta elección es natural considerar $I'=\alpha^{-1}(V') \subset \alpha^{-1}(V) \subset I \subset \bb{R}$. Por ser $\alpha$ continua tendremos que $I'$ es un abierto de $\bb{R}$ y por ser homeomorfismo $\alpha^{-1}$ será también continua lo que nos garantiza que $I'$ es conexo. Concluimos que $I'$ es un intervalo abierto de $\bb{R}$.\\

            Finalmente consideramos $\alpha' = \alpha_{|I'}$ que será una curva regular de $I'$ en $\bb{R}^2$ (por ser restricción en el dominio de una curva regular). \\
            
            Queremos ver ahora que se verifican las condiciones para que $C'$ sea una curva simple. Tenemos que 
            \begin{gather*}
                \alpha'(I') = \alpha' (\alpha^{-1} (V')) = V'
            \end{gather*}
            y además $\alpha':I'\to V'$ es un homeomorfismo por ser restricción de $\alpha$ que ya lo era.
        \end{proof}

        \item Consideramos $C\subset \bb{R}^2$ un subconjunto y $\{C_i\}_{i\in I}$ una familia de subconjuntos de $C$, es decir, $C_i\subset C$ para todo $i\in I$ tal que cada $C_i$ es una curva simple abierta en $C$ y tal que 
        \begin{gather*}
            C = \bigcup_{i\in I} C_i
        \end{gather*}
        Entonces $C$ es una curva simple.
        \begin{proof}
            La demostración se deja como ejercicio
        \end{proof}

        \item Grafos de funciones derivables. Consideramos un intervalo abierto $I\subset \bb{R}$ y una función $f: I\to \bb{R}$ es derivable. De esta forma podemos tomar
        \begin{gather*}
            Grafo(f) = \{(x,y) \in \bb{R}^2 \ : \ x\in I,\ y=f(x)\} \subset \bb{R}^2
        \end{gather*}
        y dicho grafo será una curva simple.
        \begin{proof}
            Comencemos nombrando $C=Grafo(f)\neq \emptyset$ y podemos reescribir
            \begin{gather*}
                C = \{(x,f(x)) : x\in I\}
            \end{gather*}
            Tomamos un $p=(x,f(x))\in C$ donde $x\in I$. Elegimos $V=C$ por lo que trivialmente $V$ será un entorno abierto de $p$ en $C$. También tomamos el $I$ dominio de la función que ya es un intervalo abierto de $\bb{R}$. Nos queda por escoger un $\alpha:I \to \bb{R}^2$ que a poco que pensemos es claro que nos vale
            \begin{gather*}
                \alpha(t) = (t,f(t))
            \end{gather*}  
            que es derivable por serlo la identidad y $f$. Además, tenemos que $\alpha'(t) = (1, f'(t)) \neq (0,0)$ por lo que claramente $\alpha$ es una curva regular. Tenemos que
            \begin{gather*}
                \alpha(I) = \{(t, f(t)) : t\in I\} = Grafo(f) = C = V
            \end{gather*}
            Veamos que $\alpha : I \to V$ es un homeomorfismo. Para ello consideramos $\pi:\bb{R}^2 \to \bb{R}$ la primera proyección, es decir, tal que
            \begin{gather*}
                \pi(x,y) = x
            \end{gather*}
            y se verificará que 
            \begin{gather*}
                (\pi \circ \alpha)(t) = \pi ( \alpha(t)) = \pi(t, f(t)) = t\\
                (\alpha \circ \pi)(x, f(x)) = \alpha(\pi(x, f(x))) = \alpha(x) = (x, f(x))
            \end{gather*}
            Concluimos que $\pi$ es la inversa de $\alpha$ y además es continua (ya que la proyección siempre lo es\footnote{Contenido estudiado en Topología I}) por lo que $\alpha : I \to V$ es un homeomorfismo.
        \end{proof}
        Una consecuencia clara es que todas las rectas de $\bb{R}^2$ son curvas simples. Por la primera observación tendremos también que todos los segmentos abiertos de $\bb{R}^2$ (o la unión de ellos) son curvas simples.
    \end{enumerate}
\end{ejemplo}
